\chapter{安全、权限与人在回路}
\label{chap:security}

\begin{takeaway}[学习目标]
建立 Agent 威胁模型；控制提示注入、数据泄露、越权与供应链风险；按风险设计人工确认，并理解安全规则为何必须落在模型之外。
\end{takeaway}

\section{Agent 扩大了攻击面}

普通聊天模型输出错误文字；Agent 可能读私有数据、调用 API、写文件、发消息或付款。输入来源也不再只有用户，还包括网页、邮件、文档、工具结果、记忆和第三方 Server。任何不可信内容都可能试图影响决策。

OWASP GenAI Security Project 持续整理 LLM 应用风险，是建立基线的重要入口\citep{owasp_llm_top10}。但威胁模型必须结合自己的资产、用户和工具权限。

\section{画出信任边界}

先列资产：用户数据、凭证、资金、生产环境、品牌发布权、内部知识。再列主体：用户、Host、模型提供商、MCP Server、工具后端、第三方内容和运维人员。最后画出数据跨越哪些边界，谁能发起什么动作。

每条边回答：是否认证、是否加密、是否最小权限、是否审计、是否能撤销、输入是否可信。安全评审的对象是完整数据流，不只是 system prompt。

\section{提示注入是数据与指令混淆}

直接注入来自用户；间接注入藏在网页、邮件或文档里。仅告诉模型“忽略恶意指令”不够，因为模型仍在同一上下文解释文字。

防御采用多层组合：

\begin{itemize}
  \item 明确标记外部内容为不可信数据，并隔离其边界。
  \item 检索和工具按用户权限过滤，模型不能自行扩大范围。
  \item 敏感动作由策略引擎检查，与文本内容无关。
  \item 对外发送前展示目标、内容、附件与数据来源。
  \item 高风险场景在沙箱中解析文档，禁止任意网络和代码执行。
\end{itemize}

\section{能力安全：最小权限与短期凭证}

每个工具只获得完成动作所需的 scope。只读研究 worker 不应拥有发布权限；生成报告不需要读取用户全部邮箱。凭证尽量短期、按任务签发，并绑定用户、租户、工具和资源范围。

不要把密钥放进模型上下文。工具执行器在服务端注入凭证，模型只看到能力名称与必要结果。日志也应屏蔽令牌、Cookie、个人身份信息和敏感文档。

\section{确认点按风险决定}

人在回路不是每一步都弹窗。过多确认会让用户形成机械点击。确认应该出现在不可逆、外部可见、高价值、权限扩大或低置信度的关键边界。

\begin{table}[htbp]
\centering
\begin{tabularx}{\textwidth}{>{\bfseries}p{0.19\textwidth} X X}
\toprule
风险 & 示例 & 交互 \\
\midrule
低 & 读公开资料、生成本地草稿 & 默认执行，提供活动记录 \\
中 & 修改可恢复文件、创建草稿 & 批量预览或执行后可撤销 \\
高 & 发邮件、发布内容、改共享数据 & 执行前逐项确认 \\
极高 & 付款、删生产数据、法律承诺 & 强认证、双人复核或禁止自动化 \\
\bottomrule
\end{tabularx}
\end{table}

确认令牌要绑定动作的规范化参数和版本。用户批准“发给 A 的草稿 v3”后，系统不能悄悄替换为“发给全部客户的 v4”。

\section{数据最小化与生命周期}

只收集任务所需数据；在进入模型前做字段裁剪或匿名化；区分开发日志和生产日志；为对话、trace、记忆、artifact 设置不同保留期。用户删除请求必须覆盖派生索引、缓存与备份策略。

多租户系统要在数据库查询、向量索引、缓存键、日志和 artifact 存储每一层隔离。只在 prompt 中写“不要访问其他客户”完全不构成隔离。

\section{代码执行与浏览器控制}

代码 Agent 应在临时、受限、可销毁环境中运行：非 root 用户、只挂载必要目录、网络域名白名单、CPU/内存/时间限制、禁止读取宿主密钥。提交修改前检查 diff 和测试；涉及 Git 推送、包发布或部署时独立确认。

浏览器 Agent 面对恶意页面、隐藏元素和会话 Cookie。使用独立浏览器资料、限制下载与域名、对付款或发送动作要求确认，并记录最终点击对象和页面状态。

\section{供应链与生态资源}

模型、Python 包、MCP Server、Skills、容器镜像和 Awesome 列表中的项目都属于供应链。安装前核验发布者、许可证、版本、依赖和最近维护；锁定版本，生成依赖清单，定期扫描漏洞。公开仓库的 Stars 不是安全证明。

\section{红队测试}

红队集应包含：要求泄露系统指令、网页内嵌注入、工具参数走私、跨租户查询、编码后的敏感内容、危险链接、重复批准、取消后仍执行、依赖返回恶意文本。每个样本必须检查最终业务状态，而非只看模型有没有说“拒绝”。

\begin{practice}[做一次权限演练]
给研究 Agent 加一个发送邮件工具，但默认只有草稿权限。设计 20 个攻击样本，尝试利用网页内容或用户措辞让它直接发送。验证真正的发送端点始终要求绑定具体内容与接收者的批准令牌。
\end{practice}

\begin{checkpoint}
\begin{itemize}
  \item \checkbox 我画出了资产、主体、数据流和信任边界。
  \item \checkbox 外部内容不能通过文字扩大工具权限。
  \item \checkbox 密钥不进入模型上下文，凭证按任务最小化。
  \item \checkbox 安全测试验证最终状态，而不只验证模型回答。
\end{itemize}
\end{checkpoint}

